% !TEX program = xelatex
\documentclass[a4paper,14pt,oneside,openany]{memoir}

%%% Задаем поля, отступы и межстрочный интервал %%%

\usepackage[left=30mm, right=15mm, top=20mm, bottom=20mm]{geometry}
\pagestyle{plain}
\parindent=1.25cm
\usepackage{indentfirst}

%%% Задаем языковые параметры и шрифт %%%

\usepackage[english,russian]{babel}
\babelfont[russian]{rm}{Times New Roman}
\babelfont[russian]{tt}{Courier New}

%%% Задаем стиль заголовков и подзаголовков в тексте %%%

\setsecnumdepth{subsection}
\renewcommand*{\chapterheadstart}{}
\renewcommand*{\printchaptername}{}
\renewcommand*{\chapnumfont}{\normalfont\bfseries}
\renewcommand*{\afterchapternum}{\hspace{1em}}
\renewcommand*{\printchaptertitle}{\normalfont\bfseries\centering\MakeUppercase}
\setbeforesecskip{20pt}
\setaftersecskip{20pt}
\setsecheadstyle{\raggedright\normalfont\bfseries}
\setbeforesubsecskip{20pt}
\setaftersubsecskip{20pt}
\setsubsecheadstyle{\raggedright\normalfont\bfseries}

%%% Исправление нумерации секций %%%

\renewcommand{\thesection}{\arabic{section}}
\renewcommand{\thesubsection}{\arabic{subsection}}

%%% Сквозная нумерация рисунков и таблиц %%%

\counterwithout{figure}{chapter}
\counterwithout{table}{chapter}

%%% Выравнивание и переносы %%%

\tolerance 1414
\hbadness 1414
\emergencystretch 1.5em
\clubpenalty=10000
\widowpenalty=10000

%%% Задаем параметры оформления рисунков и таблиц %%%

\usepackage{graphicx}
\usepackage{caption}
\usepackage{subcaption}
\usepackage{float}
\graphicspath{{photo/}}
\captionsetup[figure]{font=small, width=\textwidth, name=Рисунок, justification=centering}
\captiondelim{ --- }
\usepackage[section]{placeins}

%%% Задаем параметры ссылок и гиперссылок %%% 

\usepackage{hyperref}
\hypersetup{
    colorlinks=true,
    linkcolor=blue,
    citecolor=blue,
    urlcolor=blue
}

%%% Настраиваем отображение списков %%%

\usepackage{enumitem}
\renewcommand*{\labelitemi}{\normalfont{--}}
\setlist{noitemsep, leftmargin=*}

\begin{document}

% Титульный лист (без номера страницы)
\thispagestyle{empty}

\begin{center}
    МИНИСТЕРСТВО ЦИФРОВОГО РАЗВИТИЯ, СВЯЗИ И МАССОВЫХ \\ КОММУНИКАЦИЙ \\ РОССИЙСКОЙ ФЕДЕРАЦИИ

    \vspace{20pt}

    Федеральное государственное бюджетное образовательное учреждение \\ высшего образования \\
    "Сибирский государственный университет телекоммуникаций и \\ информатики"

    \vspace{20pt}

    Кафедра телекоммуникационных систем и вычислительных средств \\ (ТС и ВС)
\end{center}

\vfill

\begin{center}
    Лабораторная работа №12 \\
    по дисциплине \\
    \textit{"Web-технологии"}

    \vspace{20pt}

    по теме: \\
    \uppercase{Основы HTML и CSS}
\end{center}

\vfill

\noindent Студент: \\
\textit{Группа ИКС-532 \hfill А.С. Лёшин}

\vspace{20pt}

\noindent Преподаватель: \\
\textit{старший преподаватель \hfill А.В. Андреев}

\vfill

\begin{center}
    Новосибирск 2026 г.
\end{center}

\newpage
\setcounter{page}{2}
\OnehalfSpacing*

\subsection*{Цель работы}

\begin{itemize}
    \item Освоить базовые теги HTML для разметки текста, списков и таблиц
    \item Изучить основные CSS-селекторы и свойства оформления текста
    \item Научиться применять фоновые изображения и стилизацию списков
    \item Выполнить 16 практических задач по HTML и CSS
\end{itemize}

\section{Блок 1.1. Теги для разметки текста и атрибуты}

Выполнены задачи:
\begin{enumerate}
    \item Разметка текста с заголовком второго уровня и двумя абзацами, добавлена ссылка
    \item Исправление разметки с заголовком "Дорога", цитатой и абзацами
    \item Разметка текста с заголовками h3, фотографиями и альтернативным текстом
\end{enumerate}

\textbf{Ссылка:} \href{https://codepen.io/Dreamshelter/pen/XJjywJB}{https://codepen.io/Dreamshelter/pen/XJjywJB}

\vspace{15pt}

\section{Блок 1.2. Списки и таблицы}

Выполнены задачи:
\begin{enumerate}
    \item Маркированный список из 5 самых известных людей мира
    \item Нумерованный список ТОП-5 Forbes в порядке убывания состояния
    \item Таблица заказа аниматоров с итоговой строкой (сумма 29 350 руб.)
    \item Исправление таблицы звёзд эстрады (7 строк, 4 колонки)
\end{enumerate}

\textbf{Ссылка:} \href{https://codepen.io/Dreamshelter/pen/VYKqbeJ}{https://codepen.io/Dreamshelter/pen/VYKqbeJ}

\vspace{15pt}

\section{Блок 2.1. Селекторы и свойства (Калуша и Бутявка)}

Выполнены задачи:
\begin{enumerate}
    \item Синий цвет для заголовка "Калуша с Калушатами"
    \item Коричневый цвет для всех выделенных \texttt{strong} (Калушата)
    \item Жёлтый цвет и полужирное начертание для \texttt{span} с Бутявкой внутри цитат
    \item Курсив для \texttt{strong} внутри нумерованного списка
\end{enumerate}

\textbf{Ссылка:} \href{https://codepen.io/Dreamshelter/pen/bNwOWpB}{https://codepen.io/Dreamshelter/pen/bNwOWpB}

\vspace{15pt}

\section{Блок 2.2. Оформление текстовых блоков (Альберт Эйнштейн)}

Выполнены задачи:
\begin{enumerate}
    \item Заголовок \texttt{.main-header}: 56px, цвет \#365f91, шрифт serif
    \item Все абзацы — размер 20px
    \item Подчёркивание для \texttt{strong} внутри абзаца \texttt{.lead}
    \item Цвет списка \#365f91, маркеры: квадрат (1-й уровень), круг (2-й уровень)
    \item Цитата: фоновое изображение, цвет \#494429, курсив, шрифт serif
\end{enumerate}

\textbf{Ссылка:} \href{https://codepen.io/Dreamshelter/pen/ogzJWrb}{https://codepen.io/Dreamshelter/pen/ogzJWrb}

\section*{Заключение}
\addcontentsline{toc}{section}{Заключение}

В ходе выполнения лабораторной работы №12 были решены все задачи из четырёх блоков:

\begin{enumerate}
    \item \textbf{Блок 1.1 (Теги для разметки текста)} — 3 задачи
    \item \textbf{Блок 1.2 (Списки и таблицы)} — 4 задачи
    \item \textbf{Блок 2.1 (Селекторы и свойства)} — 4 задачи
    \item \textbf{Блок 2.2 (Оформление текстовых блоков)} — 5 задач
\end{enumerate}

\end{document}